% =============================================================================
% Module 8: Solutions to Exercises
% =============================================================================
% This file is conditionally included based on \ifsolutions
% =============================================================================

\section{Solutions to Exercises}
\label{sec:elliptical_solutions}

\noindent
Full worked solutions are also available in the Mathematica script
\solnlink{08_Elliptical_Models_Caustics/problems\_08.wl}{problems\_08.wl},
which verifies each answer symbolically and numerically.

% ----- Exercise 8.1 -----
\subsection*{Exercise~\ref{ex:sis_shear_crit}: SIS + Shear Critical Curves and Caustics}

\textbf{(a)} Starting from
$\psiL = \thetaE|\vectheta| + \frac{1}{2}\gamma_{\mathrm{ext}}
(\theta_1^2 - \theta_2^2)$:

The deflection components are:
\begin{align*}
    \alpha_1 &= \thetaE\,\frac{\theta_1}{|\vectheta|}
    + \gamma_{\mathrm{ext}}\,\theta_1 \,, \\
    \alpha_2 &= \thetaE\,\frac{\theta_2}{|\vectheta|}
    - \gamma_{\mathrm{ext}}\,\theta_2 \,.
\end{align*}

The Jacobian matrix components are:
\begin{equation*}
    \Amat_{ij} = \delta_{ij} -
    \frac{\partial\alpha_i}{\partial\theta_j} \,.
\end{equation*}
Computing the partial derivatives of the SIS deflection
($\alpha_i^{\mathrm{SIS}} = \thetaE\,\theta_i/|\vectheta|$) and
adding the shear contribution:
\begin{equation*}
    \Amat = \begin{pmatrix}
    1 - \thetaE\,\frac{\theta_2^2}{\theta^3} - \gamma_{\mathrm{ext}}
    & \thetaE\,\frac{\theta_1\theta_2}{\theta^3} \\[6pt]
    \thetaE\,\frac{\theta_1\theta_2}{\theta^3}
    & 1 - \thetaE\,\frac{\theta_1^2}{\theta^3} + \gamma_{\mathrm{ext}}
    \end{pmatrix}
\end{equation*}
where $\theta = |\vectheta|$.

The determinant is:
\begin{equation*}
    \det\Amat = 1 - \frac{\thetaE}{\theta}
    - \gamma_{\mathrm{ext}}^2
    + \frac{\thetaE}{\theta}\gamma_{\mathrm{ext}}\,
    \frac{\theta_1^2 - \theta_2^2}{\theta^2} \,.
\end{equation*}

\textbf{(b)} Setting $\det\Amat = 0$ in polar coordinates
($\theta_1 = \theta\cos\phi$, $\theta_2 = \theta\sin\phi$):
\begin{equation*}
    1 - \frac{\thetaE}{\theta} - \gamma_{\mathrm{ext}}^2
    + \frac{\thetaE}{\theta}\gamma_{\mathrm{ext}}\cos 2\phi = 0 \,.
\end{equation*}
Solving for $\theta$:
\begin{equation*}
    \theta_{\mathrm{crit}}(\phi) = \frac{\thetaE\,
    (1 - \gamma_{\mathrm{ext}}\cos 2\phi)}{1 - \gamma_{\mathrm{ext}}^2} \,.
\end{equation*}
For $\thetaE = 1''$, $\gamma_{\mathrm{ext}} = 0.1$: at $\phi = 0$,
$\theta = 0.9/0.99 \approx 0.909''$; at $\phi = \pi/2$,
$\theta = 1.1/0.99 \approx 1.111''$.  The critical curve is an
ellipse elongated along $\theta_2$. \checkmark

\textbf{(c)} Mapping through the lens equation
$\beta_i = \theta_i - \alpha_i$:
\begin{align*}
    \beta_1(\phi) &= -\gamma_{\mathrm{ext}}\,
    \theta_{\mathrm{crit}}(\phi)\cos\phi \,, \\
    \beta_2(\phi) &= \gamma_{\mathrm{ext}}\,
    \theta_{\mathrm{crit}}(\phi)\sin\phi \,.
\end{align*}
This traces out a diamond (astroid) with four cusps at $\phi = 0,
\pi/2, \pi, 3\pi/2$.  The cusps are at $\beta_1 \approx \pm 0.092''$
and $\beta_2 \approx \pm 0.112''$. \checkmark

\textbf{(d)} The area enclosed by the caustic can be computed using
the shoelace formula on the parametric curve.  For
$\gamma_{\mathrm{ext}} = 0.1$, the caustic area is approximately
$0.013\,\mathrm{arcsec}^2$.  This is verified numerically in
\texttt{problems\_08.wl}.

% ----- Exercise 8.2 -----
\subsection*{Exercise~\ref{ex:sie_caustic}: SIE Tangential Caustic}

\textbf{(a)} The SIE deflection angle components are:
\begin{align*}
    \alpha_1 &= \frac{\thetaE\sqrt{q}}{\sqrt{1-q^2}}\,
    \arctan\!\left(\frac{\theta_1\sqrt{1-q^2}}
    {\sqrt{q^2\theta_1^2 + \theta_2^2}}\right) , \\
    \alpha_2 &= \frac{\thetaE\sqrt{q}}{\sqrt{1-q^2}}\,
    \mathrm{arctanh}\!\left(\frac{\theta_2\sqrt{1-q^2}}
    {\sqrt{q^2\theta_1^2 + \theta_2^2}}\right) .
\end{align*}
For $q = 0.7$: $\sqrt{1-q^2} = \sqrt{0.51} \approx 0.714$,
$\sqrt{q} \approx 0.837$.

\textbf{(b)} The critical curve and caustic are computed numerically
by evaluating $\det\Amat$ on a fine grid (or parametrically) and
mapping the critical curve through the lens equation.  The caustic
has the diamond shape with four cusps.  The diamond is elongated
along the $\beta_2$-axis (perpendicular to the mass elongation
along $\theta_1$).

\textbf{(c)} The caustic area for $q = 0.7$ is approximately
$0.032\,\mathrm{arcsec}^2$.  This is somewhat larger than the
SIS + shear model with an ``equivalent'' shear
$\gamma_{\mathrm{ext}} \approx (1-q)/(1+q) \approx 0.176$.  The
exact comparison depends on how one maps between ellipticity and
shear; see \texttt{problems\_08.wl} for numerical details.

% ----- Exercise 8.3 -----
\subsection*{Exercise~\ref{ex:fold_magnification}: Fold Magnification Scaling}

\textbf{(a)} Near a fold caustic, we align coordinates so that the
fold is along the $\theta_\parallel$ direction and the critical
curve crosses at $\theta_\perp = 0$.  The lens mapping perpendicular
to the fold takes the form:
\begin{equation*}
    \beta_\perp = a\,\theta_\perp^2
\end{equation*}
where $a$ is a nonzero constant (this is the leading-order Taylor
expansion of the lens mapping near the fold, where the linear term
vanishes because $\det\Amat = 0$).

For $\beta_\perp > 0$ (source on the multi-image side):
$\theta_\perp = \pm\sqrt{\beta_\perp/a}$.  Two solutions exist, one
on each side of the critical curve.

For $\beta_\perp < 0$ (source on the single-image side): no real
solutions.  The images have merged and disappeared. \checkmark

\textbf{(b)} The Jacobian of this simplified mapping is:
\begin{equation*}
    \frac{d\beta_\perp}{d\theta_\perp} = 2a\,\theta_\perp \,.
\end{equation*}
The magnification in the perpendicular direction is:
\begin{equation*}
    \mu_\perp = \frac{1}{|2a\,\theta_\perp|}
    = \frac{1}{2a\sqrt{\beta_\perp/a}}
    = \frac{1}{2\sqrt{a\,\beta_\perp}} \,.
\end{equation*}
Therefore $|\mu| \propto 1/\sqrt{\beta_\perp}$, which is
$1/\sqrt{d_\perp}$ where $d_\perp$ is the distance from the
caustic. \checkmark

\textbf{(c)} The two images are at
$\theta_\perp = +\sqrt{\beta_\perp/a}$ and
$\theta_\perp = -\sqrt{\beta_\perp/a}$.  The sign of the Jacobian
determinant is:
\begin{equation*}
    \mathrm{sgn}(\det\Amat) \propto
    \mathrm{sgn}(d\beta_\perp/d\theta_\perp)
    = \mathrm{sgn}(2a\,\theta_\perp) \,.
\end{equation*}
Since the two images have $\theta_\perp$ values with opposite signs,
they have opposite signs of $\det\Amat$ and hence opposite
parities. \checkmark

% ----- Exercise 8.4 -----
\subsection*{Exercise~\ref{ex:sie_images}: SIE Image Positions}

\textbf{(a)} For $\vecbeta = (0.1, 0.2)''$ with $\thetaE = 1''$,
$q = 0.8$: The tangential caustic of this SIE has cusps at
approximately $\beta \approx 0.1''$.  The source at
$|\vecbeta| \approx 0.224''$ is outside the tangential caustic,
so we expect a \textbf{double} (2 bright images).

\textbf{(b)} Solving the lens equation numerically (see
\texttt{problems\_08.wl}), two bright images are found:
\begin{itemize}
    \item Image A: $\vectheta_A \approx (0.53, 0.79)''$
    \item Image B: $\vectheta_B \approx (-0.37, -0.52)''$
\end{itemize}
(Exact values depend on the numerical solver and are verified in
the Mathematica script.)

\textbf{(c)} The magnifications are computed from
$\mu = 1/\det\Amat$ at each image position:
\begin{itemize}
    \item $\mu_A > 0$ (Type~I, positive parity)
    \item $\mu_B < 0$ (Type~II, negative parity)
\end{itemize}
Exact values are in \texttt{problems\_08.wl}.

\textbf{(d)} For the singular SIE, the signed magnification sum
$\mu_A + \mu_B$ may be slightly negative because the ``missing''
central image (swallowed by the singularity) would contribute a
small positive magnification.  For a non-singular lens, the sum
over \textit{all} images (including the central one) satisfies
$\sum_i \mu_i > 0$.  In this example, the sum of the two
observable images is close to zero, consistent with the central
image having small positive $|\mu|$. \checkmark

% ----- Exercise 8.5 -----
\subsection*{Exercise~\ref{ex:burke_verification}: Burke's Odd-Number Theorem for the SIE}

\textbf{(a)} Source at $\vecbeta = (0.5, 0.0)''$ (outside tangential
caustic): Two images are found numerically.  For the singular SIE,
the central image is formally at the singularity and has zero flux.
In the distributional sense, there are 2 observable images (even
number).  For a regularized (non-singular) version, a faint third
image would appear near the center, giving 3 (odd). \checkmark

\textbf{(b)} Source at $\vecbeta = (0.05, 0.05)''$ (inside tangential
caustic): Four bright images are found.  Again, for the singular SIE,
the central image is lost at the singularity, giving 4 (even).  For
a non-singular version, 5 images (odd). \checkmark

\textbf{(c)} For the quad case: 2 Type~I images and 2 Type~II images
(among the 4 bright images).  Adding the central Type~III image
(for the non-singular case): 3 positive-parity images (2 Type~I +
1 Type~III) and 2 negative-parity images (2 Type~II).  The count
of positive-parity minus negative-parity images is $3 - 2 = 1$.
\checkmark

\textbf{(d)} The SIE is singular at the origin, which means the lens
potential diverges and the central image is infinitely demagnified
(effectively ``swallowed'' by the singularity).  Burke's theorem
strictly requires a smooth, non-singular lens.  For the SIE, the
theorem holds in the limiting sense: if we soften the core by any
finite amount, a faint central image appears and the total count
becomes odd.  The singular SIE is the limit of vanishing core
radius, where this central image has zero flux.  In practice, real
galaxies have finite cores, so Burke's theorem does apply, but the
central image is typically too faint to observe (demagnified by
factors of $10^3$--$10^5$).
